% =====================================================================
% INFO-H-419 Data Warehouses -- Part II: TPC-DI on ClickHouse
% Report skeleton for OpenAI Prism. Same preamble and conventions as the
% Part I report.
%
% Report editor: Bao Anh Tran (role 2). The editor writes the
% Introduction and the Conclusion, merges the sections, harmonises the
% style, exports the PDF and submits.
%
% Who writes what (docs/ROADMAP.md, section 6):
%   1    Introduction ........................... editor (Bao Anh)
%   2    TPC-DI Overview ........................ Anabelle (role 3)
%   3    Architecture ........................... Cyril (role 1)
%   4.1  Staging, Loaders and Orchestration ..... Cyril (role 1)
%   4.2  DW Schema, SCD2 and Surrogate Keys ..... Anabelle (role 3)
%   4.3  Transformation Rules ................... owner of each cluster
%   5    Validation ............................. Paula (role 4)
%   6    Benchmark Setup and Methodology ........ Bao Anh (role 2)
%   7    Results ................................ Paula (charts, analysis),
%                                                 Bao Anh (metric, 7.4)
%   8    Discussion ............................. Anabelle and Paula
%   9    Conclusion ............................. editor (Bao Anh)
%   A    Reproduction Guide ..................... Cyril (role 1)
%   B    Changes to the TPC-DI SQL .............. Paula (role 4)
%
% Conventions (as in Part I):
%   - \todo{...} marks a part still to write; none may remain in the
%     final PDF. Keep its label short: \colorbox does not break lines.
%     The guidance for each part is in the comments below its \todo.
%   - Scale factors: SF-10. Specification: Clause~7 for a whole clause,
%     \S7.5 for a sub-clause (TPC-DI v1.1.0).
%   - Code, tables and columns: \code{DimTrade}; files and settings: \path{...}.
%   - Say that our results are derived from TPC-DI, never that they are
%     TPC-DI results (TPC fair use policy).
% =====================================================================
\documentclass[11pt,a4paper]{article}

% ---------- Page layout and fonts ----------
\usepackage[a4paper, margin=2cm]{geometry}   % reduced margins
\usepackage[utf8]{inputenc}
\usepackage[T1]{fontenc}
\usepackage{lmodern}
\usepackage[english]{babel}
\usepackage{microtype}                       % better line breaks
\usepackage{parskip}                         % blank line between paragraphs, no indent

% ---------- Figures and tables ----------
\usepackage{graphicx}                        % \includegraphics (Paula's charts)
\usepackage{float}                           % [H] to place a figure exactly here
\usepackage{caption}
\usepackage{subcaption}                      % several charts side by side
\usepackage{booktabs}                        % clean tables: \toprule \midrule \bottomrule
\usepackage{tabularx}                        % tables as wide as the text
\usepackage{longtable}                       % tables over several pages
\usepackage{array}
\usepackage{multirow}

% ---------- Maths and numbers ----------
\usepackage{amsmath}                         % metric formula (Section 6.5)

% ---------- Lists, colours, code ----------
\usepackage{enumitem}
\setlist{nosep}                              % compact lists
\usepackage[table]{xcolor}
\usepackage{listings}                        % command and SQL blocks
\lstset{
  basicstyle=\ttfamily\small,
  backgroundcolor=\color{gray!10},
  frame=single, rulecolor=\color{gray!40},
  breaklines=true, breakatwhitespace=true,
  columns=fullflexible, keepspaces=true,
  showstringspaces=false,
  xleftmargin=3pt, xrightmargin=3pt,
  literate={…}{{\ldots}}1 {→}{{$\rightarrow$}}1 {×}{{$\times$}}1
           {–}{{--}}1 {—}{{---}}1
}
\lstdefinestyle{bash}{language=bash}
\lstdefinestyle{sql}{language=SQL, morekeywords={ENGINE,MergeTree,ReplacingMergeTree,OPTIMIZE,FINAL,Nullable,EXCHANGE,TABLES}}

% ---------- Links (load last) ----------
\usepackage{url}
\usepackage[hidelinks]{hyperref}             % clickable references and table of contents

% ---------- Shortcuts ----------
\newcommand{\code}[1]{\texttt{#1}}           % inline code: \code{DimTrade}
\newcommand{\todo}[1]{\colorbox{yellow}{\textbf{[TO COMPLETE: #1]}}}

\newcolumntype{L}{>{\raggedright\arraybackslash}X}   % X column, left-aligned (no stretched spaces)
\renewcommand{\arraystretch}{1.25}                    % a little more space between table rows

% ---------- Title ----------
\title{\textbf{INFO-H-419 Data Warehouses}\\[4pt]
       \Large Part II: TPC-DI on ClickHouse}
\author{Bao Anh Tran, Anabelle Meli, Cyril Stembert, Paula Ventura\\[4pt]
        \small Université libre de Bruxelles --- Lecturer: Prof.\ Esteban Zimányi}
\date{December 2026}

\begin{document}

\maketitle
\tableofcontents
\newpage

% =====================================================================
\section{Introduction}
\label{sec:intro}
% Writer: report editor (Bao Anh). Written last.

\todo{introduction, written last}
% - Goal of Part II: build an ETL that loads the TPC-DI data warehouse into
%   ClickHouse (historical load and two incremental batches) and measure it.
% - What we built, in two sentences: ELT inside ClickHouse, driven by a
%   Python orchestrator (Section 3).
% - The main results, in two sentences.
% - Structure of the report.
% - Keep this sentence (TPC fair use policy):

Our results are derived from TPC-DI and are not official TPC benchmark results~\cite{r1}.

% =====================================================================
\section{TPC-DI Overview}
\label{sec:tpcdi}
% Writer: Anabelle (role 3).

\subsection{Business Scenario}
\todo{brokerage firm and its source systems}
% The brokerage firm and its sources: OLTP database, customer management
% system (XML extract), financial newswire (FINWIRE), HR file, prospect
% list (Clause~1)~\cite{r1,r2}.

\subsection{Source Data and Staging Area}
\todo{file formats, batches, CDC}
% File formats: CSV, pipe-delimited, XML, fixed-width FINWIRE records.
% The three batches (Batch1 historical, Batch2 and Batch3 incremental),
% the CDC_FLAG / CDC_DSN fields of the incremental files (\S2.2--2.3).

\subsection{Data Warehouse Model}
\todo{dimensions, facts, history tracking}
% Dimensions, facts and reference tables; history-tracking (type 2)
% dimensions (\S4.4.1); the DImessages and Audit tables (Clause~3).

\subsection{Benchmark Phases and Metric}
\todo{phases of a benchmark run}
% Initialization, Historical Load, Incremental Updates 1 and 2, Automated
% Audit (\S7.2), in a few lines. The metric is defined in
% Section~\ref{sec:metric}.

% =====================================================================
\section{Architecture}
\label{sec:architecture}
% Writer: Cyril (role 1).

\todo{architecture diagram}
% Staging Area files -> loaders -> ClickHouse staging tables -> SQL
% transformations -> data warehouse tables, driven by the Python
% orchestrator.

\todo{tools and versions}
% Docker, ClickHouse 26.8.24.9, Python, DIGen 1.1.0 on Java 8.

\todo{why ELT inside ClickHouse}
% Decided on 6 November. Compare with dbt (dbt-clickhouse) and with a
% Python (pandas) ETL: single engine, easy timing of each step, more to
% say about ClickHouse.

% =====================================================================
\section{Implementation}
\label{sec:implementation}

\subsection{Staging, Loaders and Orchestration}
\label{sec:staging}
% Writer: Cyril (role 1).
\todo{loaders and orchestrator}
% One loader per file format; the XML conversion of CustomerMgmt.xml;
% FINWIRE records read as whole lines and cut with substring; how the
% orchestrator runs the steps in dependency order, logs them and writes
% the 'Status' message at the start of each phase.

\subsection{Data Warehouse Schema, SCD2 and Surrogate Keys}
\label{sec:schema}
% Writer: Anabelle (role 3).
\todo{table engines, sorting keys and data types}
% No column added or removed (\S3.4.5). MessageDateAndTime as
% DateTime64(3), because CURRENT_TIMESTAMP only has 1-second precision.

\todo{SCD2 and surrogate keys}
% Window functions in the historical load; rebuild + EXCHANGE TABLES in the
% incremental batches; DimTrade as a ReplacingMergeTree without a version
% column, read with FINAL. Surrogate keys without SERIAL/IDENTITY.

\subsection{Transformation Rules}
\label{sec:transformations}

\subsubsection{Cluster A: Reference Tables, DimBroker, DimCompany, DimSecurity, Financial}
% Writer: Cyril (role 1).
\todo{cluster A}
% Source to target mapping and the rules that needed care.

\subsubsection{Cluster B: DimCustomer, DimAccount, Prospect}
% Writer: Anabelle (role 3).
\todo{cluster B}
% Source to target mapping and the rules that needed care.

\subsubsection{Cluster C: DimTrade, FactCashBalances, FactHoldings}
% Writer: Bao Anh (role 2). Notes: docs/tpcdi-cluster-c.md.
\todo{cluster C}
% Status history of a trade; dimension keys taken when a trade first
% appears; running cash balance per account; holdings from DimTrade.

\subsubsection{Cluster D: FactWatches, FactMarketHistory}
% Writer: Paula (role 4).
\todo{cluster D}
% Source to target mapping and the rules that needed care.

% =====================================================================
\section{Validation}
\label{sec:validation}
% Writer: Paula (role 4).

\subsection{Batch Validation and DImessages}
\todo{batch validation, alerts and status messages}
% The batch validation query at the end of each phase (\S7.4) and the
% phase completion record; alerts and status messages.

\subsection{Data Visibility}
\todo{data visibility queries}
% The data visibility queries during the incremental phases (\S7.3), what
% the audit checks on them, and why we read with final = 1.

\subsection{Audit Results}
\todo{audit results per SF}
% Output of the audit query for each SF: number of tests, all "OK" or the
% failing tests with their explanation.

% =====================================================================
\section{Benchmark Setup and Methodology}
\label{sec:methodology}
% Writer: Bao Anh (role 2). Source: docs/benchmark-protocol-part2.md~\cite{r3}.

\subsection{System Under Test}
\todo{machine, versions, clock correction}
% Same machine and conditions as in Part I. Every clock of the WSL2 VM
% runs about 8% fast (all four clocksources tried); official phase times
% come from ClickHouse timestamps, step times are corrected by the drift
% factor measured at each run.

\subsection{Scale Factors}
\todo{reference SF and series}
% Reference SF and series; size and generation time of each SF.

\subsection{Benchmark Run}
\todo{phases and clean start}
% The phases in order, the clean start before each run, what is timed.

\subsection{Timing}
\todo{phase completion records and step timing}
% Phase completion records in DImessages, millisecond timestamps, step
% timing by the orchestrator.

\subsection{Metric}
\label{sec:metric}

The phase elapsed times come from the phase completion records:
$E_H = \mathit{CT}_1 - \mathit{CT}_0$, $E_{I1} = \mathit{CT}_2 - \mathit{CT}_1$ and
$E_{I2} = \mathit{CT}_3 - \mathit{CT}_2$ (\S7.5). With $R_H$, $R_{I1}$ and $R_{I2}$
the row counts of the three batches reported by DIGen, the metric is (\S7.6.1)
\[
  \mathit{TPC\_DI\_RPS} =
  \left\lfloor \sqrt{T_H \cdot \min(T_{I1}, T_{I2})} \right\rfloor,
  \qquad
  T_H = \frac{R_H}{E_H},
  \qquad
  T_{Ik} = \frac{R_{Ik}}{\max(E_{Ik},\, 1800)} .
\]

\todo{effect of the 1800 s floor}
% At our SFs the incremental phases are far below 1800 s, so the metric
% mostly follows the SF; we also report the unfloored throughputs.

\subsection{Repetitions and Controls}
\todo{three runs per SF}
% Three runs per SF, median with min and max; no warm runs (\S7.1).

\subsection{Deviations from the Specification}
\todo{list of deviations}
% From the protocol, section 7.

% =====================================================================
\section{Results}
\label{sec:results}
% Writers: Paula (role 4) for the charts and the analysis,
% Bao Anh (role 2) for the metric (Section 7.4).

\subsection{Phase Times per Scale Factor}
\todo{chart and table of the phase times}
% E_H, E_I1, E_I2 per SF.

\subsection{Time per Step}
\todo{which steps dominate}

\subsection{Scaling}
\todo{growth of the phase times with the SF}

\subsection{Throughput and Metric}
\label{sec:metricresults}
\todo{throughputs and TPC\_DI\_RPS per SF}
% T_H, T_I1, T_I2 and TPC_DI_RPS per SF, with the unfloored throughputs.

\subsection{Failures and Anomalies}
\todo{failures and anomalies}

% =====================================================================
\section{Discussion}
\label{sec:discussion}
% Writers: Anabelle (role 3) and Paula (role 4).

\todo{is ClickHouse a good fit for ETL?}
% UPDATE and SCD2, XML and fixed-width sources, data visibility without
% transactions. Link to the course chapters on dimension changes (dw04)
% and ETL (dw05)~\cite{r6}.

% =====================================================================
\section{Conclusion}
\label{sec:conclusion}
% Writer: report editor (Bao Anh). Written last.

\todo{conclusion, written last}
% What worked, what did not, what we would do differently.

% =====================================================================
\begin{thebibliography}{9}

\bibitem{r1}
Transaction Processing Performance Council.
\emph{TPC Benchmark DI (TPC-DI) Standard Specification}, revision 1.1.0, November 2014.
\url{https://www.tpc.org/TPC_Documents_Current_Versions/pdf/TPC-DI_v1.1.0.pdf}

\bibitem{r2}
M.~Poess, T.~Rabl, H.-A.~Jacobsen and B.~Caufield.
TPC-DI: the first industry benchmark for data integration.
\emph{Proceedings of the VLDB Endowment}, 7(13):1367--1378, 2014.

\bibitem{r3}
Benchmark protocol, Part II. \path{docs/benchmark-protocol-part2.md} in~\cite{r7}.

\bibitem{r4}
TPC-DI audit, batch validation and data visibility queries (Appendices A--C of the
specification), copy in the \code{databricks-tpc-di} repository,
folder \path{src/incremental_batches/audit_validation/Appendix_v1.1.0}.
\url{https://github.com/shannon-barrow/databricks-tpc-di}

\bibitem{r5}
ClickHouse documentation. \url{https://clickhouse.com/docs}

\bibitem{r6}
E.~Zimányi. INFO-H-419 Data Warehouses, course slides (dw04 Dimension Changes,
dw05 ETL). Université libre de Bruxelles.
\url{https://cs.ulb.ac.be/public/teaching/infoh419}

\bibitem{r7}
Project repository: scripts, DDL, ETL, results and analysis.
\url{https://github.com/cyyr1l/info-h419-clickhouse}

\end{thebibliography}

% =====================================================================
\appendix

\section{Reproduction Guide}
\label{app:reproduction}
% Writer: Cyril (role 1).
\todo{reproduction guide}
% Requirements, the one-command run, expected time and disk space per SF.

\section{Changes to the TPC-DI SQL}
\label{app:sql}
% Writer: Paula (role 4).
\todo{changes to the audit, validation and visibility SQL}
% Every change made to the queries of Appendices A--C to run them on
% ClickHouse, with its reason (\S3.4.11)~\cite{r4}.

\end{document}
